\subsubsection{Equações de Lagrange e estrutura das forças de inércia}
\label{sec:dinAcoplada_equacoes_Lagrange}

As equações de Lagrange para o sistema base--\gls{mr}, em termos das coordenadas generalizadas $\vec q$ e de suas derivadas, podem ser escritas componente a componente como
\cite{Wilde2018}

\begin{equation}
\frac{d}{dt}\!\left(\frac{\partial T}{\partial \dot q_k}\right)
-\frac{\partial T}{\partial q_k}
=
\tau_k,
\qquad
k = 1,\dots,n_g,
\label{eq:lagrange_componentes}
\end{equation}

onde $T$ é a energia cinética total, $\tau_k$ é a força generalizada associada à coordenada $q_k$ e $n_g$ é o número de graus de liberdade considerados na dinâmica acoplada (base e juntas).

Admitindo que a energia cinética possa ser escrita de forma quadrática em função das velocidades
generalizadas,

\begin{equation}
T(\vec q,\dot{\vec q})
=
\frac{1}{2}\,\dot{\vec q}^{\mathsf T} M(\vec q)\,\dot{\vec q},
\label{eq:T_quadratica}
\end{equation}

os elementos $m_{kj}$ da matriz de inércia generalizada $M(\vec q)$ são definidos por

\begin{equation}
T =
\tfrac{1}{2}\sum_{k=1}^{n_g}\sum_{j=1}^{n_g} m_{kj}(\vec q)\,\dot q_k \dot q_j.
\end{equation}

A partir dessa forma, \cite{Geradin2004} mostra que o primeiro termo de
\eqref{eq:lagrange_componentes} pode ser reescrito como

\begin{equation}
\frac{d}{dt}\!\left(\frac{\partial T}{\partial \dot q_k}\right)
=
\sum_{j=1}^{n_g} m_{kj}\,\ddot q_j
+
\sum_{j=1}^{n_g}\sum_{l=1}^{n_g}
\frac{\partial m_{kj}}{\partial q_l}\,\dot q_l \dot q_j,
\label{eq:lagrange_primeiro_termo}
\end{equation}

enquanto o segundo termo assume a forma

\begin{equation}
\frac{\partial T}{\partial q_k}
=
\frac{1}{2}
\sum_{j=1}^{n_g}\sum_{l=1}^{n_g}
\frac{\partial m_{jl}}{\partial q_k}\,\dot q_j \dot q_l.
\label{eq:lagrange_segundo_termo}
\end{equation}

Combinando \eqref{eq:lagrange_primeiro_termo} e \eqref{eq:lagrange_segundo_termo},
o resultado de \eqref{eq:lagrange_componentes} pode ser organizado em forma matricial
\cite{Geradin2004}:

\begin{equation}
M(\vec q)\,\ddot{\vec q}
+
C(\vec q,\dot{\vec q})\,\dot{\vec q}
=
\vec\tau
-
\vec G(\vec q),
\label{eq:dinamica_geral_com_G}
\end{equation}

em que o segundo termo representa as contribuições centrífugas e de Coriolis, e a matriz
$C(\vec q,\dot{\vec q})$ possui elementos

\begin{equation}
c_{kj}(\vec q,\dot{\vec q})
=
\sum_{l=1}^{n_g}
\bigg(
\frac{\partial m_{kj}}{\partial q_l}
-
\frac{1}{2}\frac{\partial m_{lj}}{\partial q_k}
\bigg)\,\dot q_l,
\label{eq:coeficientes_C_geradin}
\end{equation}

enquanto $\vec G(\vec q)$ agrupa as forças generalizadas devidas ao potencial gravitacional
\cite{Geradin2004}.
Uma propriedade importante é que a matriz

\begin{equation}
S(\vec q,\dot{\vec q})
=
\dot M(\vec q) - 2\,C(\vec q,\dot{\vec q}),
\label{eq:S_Mdot_2C}
\end{equation}

obtida a partir dos elementos $m_{kj}$, é antissimétrica, ou seja, possui elementos $s_{kj} = -\,s_{jk}$ e satisfaz $S = -S^{\mathsf T}$.

No contexto deste trabalho, em regime de microgravidade, o potencial gravitacional é omitido e o vetor $\vec G(\vec q)$ é nulo.
Portanto, a equação de movimento do conjunto base--\gls{mr} é finalmente escrita como

\begin{equation}
M(\vec q)\,\ddot{\vec q}
+
C(\vec q,\dot{\vec q})\,\dot{\vec q}
=
\vec\tau,
\label{eq:dinamica_geral}
\end{equation}

em que $M(\vec q)$ é a matriz de inércia generalizada obtida a partir das velocidades espaciais da base e dos elos, e $C(\vec q,\dot{\vec q})\,\dot{\vec q}$ representa o vetor de forças centrífugas--Coriolis \cite{lynch2017modern}.
%%%%%%%%%%%%%%%%%%%%%%%%%%%%%%%%%%%%%%%%%%%%%%%%%%%%%%%%%%%%%
%%%%%%%%%%%%%%%%%%%%%%%%%%%%%%%%%%%%%%%%%%%%%%%%%%%%%%%%%%%%%
%%%%%%%%%%%%%%%%%%%%%%%%%%%%%%%%%%%%%%%%%%%%%%%%%%%%%%%%%%%%%
%%%%%%%%%%%%%%%%%%%%%%%%%%%%%%%%%%%%%%%%%%%%%%%%%%%%%%%%%%%%%
Sendo assim, é conveniente explicitar a Equação~\eqref{eq:dinamica_geral} entre base--\gls{mr} por meio do particionamento, conforme \cite{featherstone2008,Wilde2018} da matriz de inércia e da matriz de Coriolis:

\begin{equation}
M(\vec q) =
\begin{bmatrix}
M_{BB}(\vec q) & M_{BM}(\vec q) \\
M_{MB}(\vec q) & M_{MM}(\vec q)
\end{bmatrix},
\qquad
C(\vec q,\dot{\vec q}) =
\begin{bmatrix}
C_{BB}(\vec q,\dot{\vec q}) & C_{BM}(\vec q,\dot{\vec q}) \\
C_{MB}(\vec q,\dot{\vec q}) & C_{MM}(\vec q,\dot{\vec q})
\end{bmatrix},
\label{eq:particionamento_M_C}
\end{equation}

onde $B$ são as contribuições da inércia da base física e $M$ são as contribuições da inércia do \gls{mr}. Então, obtém-se o seguinte formato

\begin{equation}
M_{BB}(\vec q)\,\ddot{\vec q}_B
+
M_{B\theta}(\vec q)\,\ddot{\vec\theta}
+
C_{BB}(\vec q,\dot{\vec q})\,\dot{\vec q}_B
+
C_{B\theta}(\vec q,\dot{\vec q})\,\dot{\vec\theta}
=
\vec\tau_{reacao},
\label{eq:dinAcoplada_base_compacta}
\end{equation}

onde $\tau_{reacao}$ é a perturbação exercida na base física, causada pela movimentação do \gls{mr}.
